%% Chapter 4: Results
\chapter{Findings and Results}
\label{ch:results}

This chapter presents the experimental findings organised around the research objectives defined in Chapter~\ref{ch:introduction}. Section~\ref{sec:res-main} reports the core classification results across the four-condition evaluation matrix. Section~\ref{sec:res-features} examines the effect of feature reduction. Section~\ref{sec:res-ablation} presents the noise component ablation. Section~\ref{sec:res-robustness} evaluates robustness under severity scaling, and Section~\ref{sec:res-extended} summarises extended analyses including calibration, domain gap quantification, and per-class vulnerability profiling. All results report mean $\pm$ standard deviation over 5 random seeds (independent stratified splits) unless otherwise stated.

%% ─────────────────────────────────────────────────────────────────────────────
\section{Main Classification Results}
\label{sec:res-main}

Table~\ref{tab:main-results-raw} presents the full results matrix for all four models on raw (208-D) features across the four evaluation conditions.

\begin{table}[htbp]
\centering
\caption{Classification accuracy (\%) on raw features (208-D) across four evaluation conditions. Mean $\pm$ std over 5 seeds. Bold indicates best model per condition.}
\label{tab:main-results-raw}
\small
\begin{tabular}{@{}lcccc@{}}
\toprule
\textbf{Model} & \textbf{Clean$\to$Clean} & \textbf{Clean$\to$Noisy} & \textbf{Noisy$\to$Noisy} & \textbf{Noisy$\to$Clean} \\
\midrule
CNN  & 90.5 $\pm$ 5.3 & 19.4 $\pm$ 0.8 & \textbf{76.1 $\pm$ 0.6} & \textbf{52.6 $\pm$ 1.1} \\
MLP  & \textbf{98.8 $\pm$ 0.7} & 21.3 $\pm$ 0.3 & 57.3 $\pm$ 0.8 & 47.6 $\pm$ 0.7 \\
RF   & 91.0 $\pm$ 0.3 & 22.6 $\pm$ 0.2 & 57.9 $\pm$ 0.6 & 29.9 $\pm$ 3.0 \\
SVM  & 83.4 $\pm$ 0.4 & 20.0 $\pm$ 0.0 & 50.9 $\pm$ 0.7 & 51.7 $\pm$ 0.4 \\
\bottomrule
\end{tabular}
\end{table}

Three findings emerge immediately. First, all models achieve strong clean-domain performance (83--99\%), confirming the forward model produces well-separated class representations amenable to both classical and deep learning approaches. Second, all models collapse to approximately chance level ($\approx$20\%) under the Clean$\to$Noisy condition demonstrating that the multi-component domain gap is a \emph{distributional} problem rather than a capacity problem. Third, noise-aware training substantially recovers performance. The CNN achieves 76.1\% on Noisy$\to$Noisy, representing a +56.7 percentage point (pp) improvement over Clean$\to$Noisy ($p < 0.0001$, Cohen's $d = -135.9$). This constitutes the primary evidence that the multi-component noise model enables effective noise-matched training, directly addressing Research Objective~3.

The CNN significantly outperforms all baselines on noisy data (+18.2pp over the next-best RF at 57.9\%), demonstrating that convolutional architectures exploit the structured, per-electrode nature of the noise through learned local filters that span adjacent measurements within the same electrode block. Conversely, on clean data the MLP achieves the highest accuracy (98.8\%), while the CNN exhibits training instability (std = 5.3\%) manifesting as oscillating validation loss across seeds---suggesting over-parameterisation for the simpler, noise-free discrimination task where all classes are linearly separable in the raw feature space. This inverted architecture ranking between clean and noisy conditions reveals that the inductive biases conferring noise robustness (local receptive fields, weight sharing) impose a small accuracy cost when noise structure is absent.

The training dynamics provide further insight into the noise-matched learning process. The Noisy$\to$Noisy CNN (Figure~\ref{fig:app-training-noisy}, Appendix~\ref{app:training-curves}) converges within approximately 60 epochs with smooth, monotonically decreasing validation loss---considerably slower than the Clean$\to$Clean model (Figure~\ref{fig:app-training-clean}), which converges in under 20 epochs. This difference reflects the increased complexity of the noise-corrupted decision boundaries. The confusion matrix for the Noisy$\to$Noisy condition (Figure~\ref{fig:app-cm-noisy}) reveals that the dominant confusions occur between light touch/point contact and between firm press/distributed contact---class pairs with similar voltage magnitude but different spatial extent, consistent with the noise model's corruption of spatial information while preserving gross magnitude.

Table~\ref{tab:augmented-results} presents the augmented and mixed training strategies for the CNN.

\begin{table}[htbp]
\centering
\caption{CNN accuracy (\%) under augmented and mixed training regimes on raw features.}
\label{tab:augmented-results}
\small
\begin{tabular}{@{}lcc@{}}
\toprule
\textbf{Training Strategy} & \textbf{Eval Clean} & \textbf{Eval Noisy} \\
\midrule
Fixed noisy   & 52.6 $\pm$ 1.1 & 76.1 $\pm$ 0.6 \\
Augmented (online noise) & 98.5 $\pm$ 0.4 & 19.8 $\pm$ 0.5 \\
Mixed (30\% clean + online noise) & 98.7 $\pm$ 0.3 & 18.2 $\pm$ 1.9 \\
\bottomrule
\end{tabular}
\end{table}

Online noise augmentation achieves near-perfect clean-domain accuracy (98.5--98.7\%) but fails completely on noisy evaluation---statistically indistinguishable from Clean$\to$Noisy ($t = 1.74$, $p = 0.78$). This reveals a fundamental paradox: random-per-epoch noise augmentation teaches the CNN to \emph{extract the persistent clean signal} rather than to classify under fixed noise. The model converges to a clean-signal extraction function that actively \emph{subtracts} the variable noise---precisely because the noise changes each epoch while the underlying signal remains constant. Training curves (Appendix~\ref{app:training-curves}, Figure~\ref{fig:app-training-augmented}) confirm stable convergence without overfitting---the model reaches low training loss and produces well-separated clean features---indicating that this is a \emph{representation learning} phenomenon rather than a training failure. The confusion matrices for augmented and clean-trained models (Appendix~\ref{app:confusion-matrices}, Figures~\ref{fig:app-cm-augmented}--\ref{fig:app-cm-clean-noisy}) exhibit near-identical failure patterns, confirming both models have converged to the same clean-signal representation. This finding has significant theoretical implications for domain randomisation \parencite{tobin2017} when applied to sensor modalities with fixed per-device characteristics, discussed in Chapter~\ref{ch:discussion}.

%% ─────────────────────────────────────────────────────────────────────────────
\section{Effect of Feature Reduction}
\label{sec:res-features}

Table~\ref{tab:feature-reduction} compares the best model per condition across the three feature representations.

\begin{table}[htbp]
\centering
\caption{Best accuracy (\%) per condition across feature representations. Feature reduction is computed on clean training data.}
\label{tab:feature-reduction}
\small
\begin{tabular}{@{}lccc@{}}
\toprule
\textbf{Condition} & \textbf{Raw (208-D)} & \textbf{PCA (7-D)} & \textbf{LDA (4-D)} \\
\midrule
Clean$\to$Clean & 98.8 (MLP) & 92.6 (MLP) & 91.4 (RF) \\
Clean$\to$Noisy & 22.6 (RF) & 24.4 (RF) & 20.0 (SVM) \\
Noisy$\to$Noisy & 76.1 (CNN) & 38.5 (SVM) & 38.3 (RF) \\
Noisy$\to$Clean & 52.6 (CNN) & 41.9 (MLP) & 22.7 (RF) \\
\bottomrule
\end{tabular}
\end{table}

Dimensionality reduction computed on clean data is counterproductive under noise: PCA/LDA reduce Noisy$\to$Noisy from 76.1\% to $\approx$38\%. The projections---computed from the covariance structure of clean data---discard precisely those dimensions that encode the noise structure a model must learn to accommodate. This empirically confirms the theoretical prediction of \textcite{ben-david2010domain}: feature representations that minimise source-domain variance may maximise target-domain error when the distributional shift introduces variance along orthogonal directions. Notably, on reduced features all models converge to $\approx$38\% regardless of architecture, suggesting the reduced spaces retain only the no-contact (class~0) discriminability---which contributes 100\%/5 = 20\% baseline even when other classes approach chance.

On clean data, LDA features preserve discriminability surprisingly well (91.4\% with only 4 dimensions), confirming the supervised projection captures the most class-relevant directions. PCA (7 components, 85.7\% variance) captures sufficient structure for clean classification but discards noise-relevant information.

%% ─────────────────────────────────────────────────────────────────────────────
\section{Noise Component Ablation}
\label{sec:res-ablation}

Table~\ref{tab:ablation-single} presents the single-component isolation results from the ablation study. Note that the ablation experiments employ a dedicated single-seed protocol with extended training budget (200 epochs, patience 30) to maximise per-condition convergence, explaining the higher absolute accuracies compared to the main 5-fold cross-validated results in Table~\ref{tab:main-results-raw}. Crucially, both training \emph{and} evaluation use full 4-component noise, so the high accuracies (97--98\%) reflect the matched-distribution ceiling for this specific training configuration rather than contradicting the main experiment's 76.1\%---which averages over 5 random seeds with different data splits and a 100-epoch budget.

\begin{table}[htbp]
\centering
\caption{CNN accuracy (\%) when trained with individual noise components and evaluated on the full noise model. All models are evaluated against the complete 4-component noise.}
\label{tab:ablation-single}
\small
\begin{tabular}{@{}lcc@{}}
\toprule
\textbf{Training Noise} & \textbf{Test Accuracy} & \textbf{$\Delta$ vs Full Model} \\
\midrule
Full 4-component & 98.3 & --- \\
Electrode bias only & 97.1 & $-$1.2 \\
Contact impedance only & 90.9 & $-$7.4 \\
Gaussian only & 92.9 & $-$5.4 \\
Quantisation only & 82.0 & $-$16.3 \\
\bottomrule
\end{tabular}
\end{table}

Quantisation produces the most harmful isolated impact (82.0\%), representing irrecoverable information loss through bit-depth reduction. However, the severity sweep (Figure~\ref{fig:severity-sweep}) reveals that within the full noise model, quantisation contributes negligibly---its effects are masked by the dominant structured components.

\begin{figure}[htbp]
\centering
\includegraphics[width=0.85\textwidth]{../results/figures/ablation/per_component_severity_sweep.png}
\caption{Per-component severity sweep. Model trained on full noise at 1$\times$ severity, evaluated with only one noise type active at varying severity. Contact impedance causes the fastest degradation ($-$4.5pp from 0$\times$ to 3$\times$).}
\label{fig:severity-sweep}
\end{figure}

The component ordering analysis reveals negligible effect size: the optimal ordering (contact impedance $\to$ Gaussian $\to$ electrode bias $\to$ quantisation) achieves 98.4\%, only 0.37pp above the worst permutation. Incrementally adding components yields diminishing returns beyond three: 97.1\% (1 component) $\to$ 98.9\% (2) $\to$ 99.3\% (3) $\to$ 99.3\% (4).

%% ─────────────────────────────────────────────────────────────────────────────
\section{Robustness Under Severity Scaling}
\label{sec:res-robustness}

\begin{figure}[htbp]
\centering
\includegraphics[width=0.85\textwidth]{../results/figures/noise_type_severity_sweep.png}
\caption{Noise-type severity sweep showing per-component degradation profiles. The full-noise-trained CNN maintains $>$94.5\% accuracy even at 3$\times$ combined severity, demonstrating robust generalisation beyond training-time noise levels.}
\label{fig:noise-type-sweep}
\end{figure}

The full-noise-trained CNN degrades gracefully from 98.9\% at 0$\times$ to 94.5\% at 3$\times$ combined severity (Figure~\ref{fig:noise-type-sweep}), demonstrating robustness well beyond training-time noise levels---a desirable property absent from prior \gls{eit} noise studies \parencite{chen2023, yao2022}. Contact impedance drives the fastest degradation ($-$1.9pp per severity unit), consistent with its multiplicative nature: scaling contact impedance variation compounds exponentially with severity, whereas additive components scale linearly. Quantisation is effectively invisible (0.0pp change), confirming that 16-bit resolution provides effectively infinite precision for the voltage magnitudes encountered. Gaussian noise causes minimal degradation ($-$0.2pp total), consistent with the finding that modern batch-normalised networks are inherently robust to additive zero-mean perturbation \parencite{rusak2020simple}---a result that further underscores why single-component Gaussian models fail to expose the real challenge.

%% ─────────────────────────────────────────────────────────────────────────────
\section{Statistical Significance and Extended Analyses}
\label{sec:res-extended}

\subsection{Statistical Testing}

Paired $t$-tests with Bonferroni correction confirm the key comparisons:

\begin{itemize}
    \item Noisy training vs clean training on noisy data: $t = -303.9$, $p < 0.0001$, Cohen's $d = -135.9$ (highly significant, very large effect).
    \item Clean$\to$Noisy vs Augmented$\to$Noisy: $t = 1.74$, $p = 0.78$ (not significant---both fail equally).
    \item Noisy$\to$Noisy vs Noisy$\to$Clean: $t = 26.5$, $p < 0.001$, $d = 11.9$ (significant asymmetry).
\end{itemize}

\subsection{Per-Class Robustness}

\begin{figure}[htbp]
\centering
\includegraphics[width=0.85\textwidth]{../results/figures/per_class_robustness.png}
\caption{Per-class accuracy as a function of noise severity. Light touch is the most vulnerable class, losing 14.7pp at 3$\times$ severity. No contact maintains 100\% throughout.}
\label{fig:per-class}
\end{figure}

Figure~\ref{fig:per-class} reveals that \emph{light touch} is the most noise-vulnerable class ($-$14.7pp at 3$\times$ severity), followed by firm press ($-$4.3pp) and distributed contact ($-$3.3pp). No contact maintains perfect classification at all severity levels, and point contact paradoxically shows slight improvement (+0.3pp). This vulnerability hierarchy correlates with signal magnitude: light touch produces the smallest voltage perturbation L2 norm (0.004), approaching the noise floor earliest as severity scales. Physically, light touch induces minimal conductivity change in the sensing medium, generating boundary voltages whose magnitude is comparable to electrode bias perturbations at elevated severity. By contrast, no-contact classification relies on detecting the \emph{absence} of any perturbation---a structurally distinct signature that noise can obscure but cannot mimic. The robustness asymmetry between no-contact (100\% at all severities) and light touch (85.3\% at 3$\times$) has direct implications for prosthetic grasp control: the most safety-critical transition---detecting initial contact onset---is precisely where the model is most susceptible to noise-induced misclassification.

\subsection{Gaussian-Only Comparison}

Under the Gaussian-only noise assumption common in prior \gls{eit} literature \parencite{chen2023, yao2022}, all models achieve $>$99\% accuracy even at 20\,dB SNR (Table~\ref{tab:gaussian-only}). The clean-trained CNN loses merely 0.2pp across a 40\,dB range---four orders of magnitude in noise power---demonstrating conclusively that Gaussian perturbation alone cannot expose the classification challenge. This result is unsurprising given the neural architecture: batch normalisation \parencite{ioffe2015batchnorm} is provably robust to additive zero-mean noise because the layer-wise standardisation absorbs Gaussian shifts by construction. The multi-component noise model developed in this work (Section~\ref{sec:method-noise}) therefore represents a necessary methodological advance: without electrode bias and contact impedance variation, any ML evaluation on simulated \gls{eit} data would report near-ceiling performance that fundamentally mischaracterises deployment readiness.

\begin{table}[htbp]
\centering
\caption{CNN accuracy (\%) under Gaussian-only noise at varying SNR levels.}
\label{tab:gaussian-only}
\small
\begin{tabular}{@{}lccccc@{}}
\toprule
\textbf{Model} & \textbf{60\,dB} & \textbf{50\,dB} & \textbf{40\,dB} & \textbf{30\,dB} & \textbf{20\,dB} \\
\midrule
CNN (clean-trained) & 99.7 & 99.7 & 99.6 & 99.6 & 99.5 \\
CNN (full-noise-aug.) & 98.9 & 98.9 & 98.9 & 98.9 & 98.9 \\
\bottomrule
\end{tabular}
\end{table}

\subsection{Confidence Calibration}

The CNN exhibits catastrophic miscalibration under domain shift: when evaluated on noisy data after clean training, it achieves ECE = 0.795 with mean confidence 99.6\% but actual accuracy of only 20.1\% \parencite{guo2017calibration}. This represents the worst possible calibration failure mode for a safety-critical system---systematic overconfidence on incorrect predictions. The reliability diagram reveals that predicted probabilities concentrate in the 0.95--1.0 bin regardless of actual correctness, indicating that the softmax output provides no discriminative uncertainty signal under domain shift.

Noise-augmented models improve clean-domain calibration substantially (ECE = 0.053), demonstrating that training on noisy data induces more diffuse softmax distributions. However, these models remain poorly calibrated on noisy data (ECE = 0.773), suggesting that even matched-distribution training does not produce well-calibrated neural networks---consistent with the general finding of \textcite{guo2017calibration} that modern architectures are overconfident by construction. This identifies post-hoc calibration (temperature scaling) or explicit uncertainty quantification (Monte Carlo dropout, deep ensembles) as prerequisites for any deployment pathway.

\subsection{Domain Gap Quantification}

Maximum Mean Discrepancy analysis \parencite{gretton2012kernel} with an RBF kernel reveals that electrode bias dominates the distributional shift (per-component MMD: electrode bias = 0.387, Gaussian = 0.013, contact impedance $\approx$ 0, quantisation $\approx$ 0). This aligns with the dataset validation finding that electrode bias contributes 99.8\% of total noise energy and provides a formal quantification of the domain adaptation challenge: the overall MMD of 0.384 represents substantial distributional divergence, and---following the theoretical bounds of \textcite{ben-david2010domain}---predicts that the target-domain error cannot be reduced below approximately $0.384 \times \lambda$ without explicit domain alignment (where $\lambda$ is a problem-dependent constant). The t-SNE visualisations (Figures~\ref{fig:app-tsne}--\ref{fig:app-tsne-components}, Appendix~\ref{app:tsne}) provide qualitative confirmation: electrode bias produces severe cluster disruption while Gaussian noise and quantisation cause minimal visible effect on class structure. The component-wise decomposition of MMD constitutes a novel diagnostic: it reveals that the ``simulation-to-reality gap'' for \gls{eit} tactile sensing is not a diffuse, multi-source problem but rather a single dominant structured distortion amenable to targeted mitigation.

\subsection{Sensitivity Analyses}

Noise parameter sensitivity testing (SNR: 20--60\,dB; contact impedance $\sigma$: 5--25\%; electrode bias: 0.005--0.08; ADC: 8--16 bits) confirms accuracy remains above 97.5\% across all physically plausible parameter ranges except extreme Gaussian noise (20\,dB: 88.9\%). Hyperparameter sensitivity (learning rate, dropout, weight decay, severity range) shows narrow clean-domain accuracy ranges ($\pm$2pp), confirming results are insensitive to specific hyperparameter choices and represent robust findings generalising beyond the specific training configuration.

\subsection{Noise Memorisation Analysis}
\label{sec:res-memorisation}

Two additional experiments directly test whether the Noisy$\to$Noisy CNN has learned genuine noise accommodation or merely memorised the training-set noise distribution.

\paragraph{Different noise draw test.} The trained noisy-domain CNN was evaluated on 10 independently-generated noisy test sets (same parametric noise distribution, different random seeds). Accuracy remained stable at $58.3\% \pm 0.7\%$ across alternative noise draws (original test set: 57.7\%), with a negligible accuracy change of $-0.6$pp (range: 57.3--59.3\%). Since each alternative test set uses entirely different electrode bias realisations, contact impedance factors, and Gaussian noise instances, this stability provides direct empirical evidence that the model generalises across noise instances rather than memorising specific bias patterns from the training data. The model has learned distributional noise accommodation---recognising the \emph{statistical structure} of multi-component noise---rather than memorising particular noise vectors.

\paragraph{Fixed-bias augmentation.} A CNN was trained with noise sampled \emph{once} per training sample and held fixed across all epochs, mimicking the deployment scenario where each physical device has persistent per-electrode characteristics. This strategy achieved 91.0\% on clean evaluation but only 19.9\% on noisy evaluation---statistically indistinguishable from both online augmentation (19.8\%) and the Clean$\to$Noisy baseline (19.4\%). The model converged in 41 epochs with a validation loss of 23.0, substantially higher than the fixed noisy dataset model (val loss $<$ 1.0), indicating failure to learn meaningful noise-domain representations.

This result refines the augmentation paradox: the critical factor is not merely noise \emph{persistence} across epochs (which fixed-bias provides) but noise \emph{consistency between train and test distributions}. In the fixed noisy dataset, the specific noise realisation applied to each sample at generation time is identical in both training and evaluation; in fixed-bias augmentation, the noise is persistent per-sample but drawn from a \emph{different realisation} than the test set. The model requires exposure to the \emph{same distributional instance} (or at minimum, noise drawn from an overlapping statistical regime) to develop effective accommodation strategies. This identifies a strict requirement for sim-to-real deployment: the training noise model must match the target hardware's noise characteristics not merely parametrically but in its per-device persistence structure.
